\section{A roadmap for groups with device and fabrication resources}\label{sec:roadmap}
The simulations narrow what has to be shown in hardware. We propose four stages. Each has pass criteria taken from this paper's
specification sheet (Tables~\ref{tab:spec} and~\ref{tab:specrel}) and main results, so a group can decide quickly whether its
device is worth an array.

\paragraph{Stage 1: single-device characterisation} (tens of devices, no array).
\begin{itemize}
\item \textbf{Measurement.} Apply at least 1000 up and 1000 down pulses at the intended pulse width and voltage, and repeat at
      least three times on at least 32 devices. Fit our device model to the records:
      \begin{itemize}
      \item step against state (giving $\gamma_\pm$);
      \item up/down asymmetry $a$;
      \item cycle-to-cycle noise $c_{2c}$ (scatter over step);
      \item device-to-device spreads;
      \item drift time constant.
      \end{itemize}
      The fitting script and the catalog format are in the release.
\item \textbf{Pass for TTv2-class training} (others at a protonic-like base):
      \begin{itemize}
      \item asymmetry \RSpecrelTtvTwoCifarWThreeTwoAsym;
      \item saturation $0<\gamma$ and $\gamma$ \RSpecrelTtvTwoCifarWThreeTwoGamma{};
      \item write noise \RSpecrelTtvTwoCifarWThreeTwoCTwoc{} steps;
      \item levels \RSpecrelTtvTwoCifarWThreeTwoNStates;
      \item retention $\tau$ \RSpecrelTtvTwoCifarWThreeTwoTauSteps{} updates.
      \end{itemize}
      These are the CIFAR $c=32$ degradation thresholds. The LM thresholds are similar (Table~\ref{tab:specrel}).
\item \textbf{Endurance.} On the Li-ECRAM and protonic devices, the tolerant algorithms applied at most \RPpwrLmMax{} pulses per weight per row on our LM. A
      training run over $10^{10}$ rows therefore needs about \REndurTenTen{} pulses per weight, shared by its two devices. The more than $10^9$ updates reported for protonic ECRAM in
2026~\cite{xu2026} would cover this.
\end{itemize}

\paragraph{Stage 2: the symmetry point.} This is the specification the simulations single out.
\begin{itemize}
\item \textbf{Measurement.} Program each device's reference to its symmetry point by alternating pulses, and measure the
      residual reference error $\sigma_r$ (in units of the half-range) and its drift over hours and over $10^6$ cycles.
\item \textbf{Pass for TTv2:} $\sigma_r$ \RSpecrelTtvTwoCifarWThreeTwoRefErr{}, stable over the training time.
\item \textbf{Otherwise:} show that chopped TTv2 or AGAD tolerates the measured offsets. Our validation reproduces their
      robustness up to $\sigma_r=0.5$ on a coarse device (Section~\ref{sec:validation}).
\end{itemize}

\paragraph{Stage 3: small arrays with hardware in the loop.}
\begin{itemize}
\item \textbf{Hardware.} Two $64\times64$ arrays per layer, the $A$ and $C$ arrays of the Tiki-Taka family.
\item \textbf{Protocol.}
      \begin{enumerate}
      \item Run the forward and backward passes digitally on the host from the arrays' read weights.
      \item Apply each update on chip as stochastic pulse coincidences on $A$, with column reads of $A$ into a host buffer $H$
            and threshold pulses into $C$.
      \item Reset or chop as the algorithm requires.
      \end{enumerate}
\item \textbf{Workloads.} A linear-regression test (a target matrix learned to below 5\% of its untrained error, as in our
      checks), then an MNIST MLP on $8\times8$ downsampled digits (64 inputs).
\item \textbf{Prediction first.} Before training, fit the Stage~1 device model, run this paper's simulator with that fit, and
      record the predicted accuracy and its seed interval. The pass criterion is that the hardware lands inside the predicted
      interval and within 2\,pp of FP. The second condition shows the device works; the first makes the simulator falsifiable.
\end{itemize}

\paragraph{Stage 4: full arrays and in-memory reads.}
\begin{itemize}
\item \textbf{Hardware.} $256\times256$ tiles, enough for the 784--256--128--10 MLP with six to eight tiles per array type,
      then the CIFAR network of this paper.
\item \textbf{Measurement.} Move the forward and backward MVMs on chip. Measure the system energy per weight per row with the
      periphery, against this paper's model (Section~\ref{sec:energy}).
\item \textbf{The energy target} is reads at about \REnNeedTops\,TOPS/W including the ADCs, for a 100$\times$ advantage over an
      H100 at its measured utilisation.
\end{itemize}

\paragraph{The 1960 memistor.} Appendix~\ref{app:bench} gives a bench replication with hobby parts and safety notes. It can
measure the memistor's missing parameters and place the cell on the same specification sheet. It is a historical and
educational experiment: by energy, the memistor cannot compete.

\paragraph{What to report.} For any device, report the fields of our catalog with their provenance: levels, $\gamma_\pm$, $a$,
$c_{2c}$, spreads, retention, reference error, pulse energy and voltage. Fitted parameters with their fitting procedure are
more useful to the field than headline level counts.

\section{Conclusion}
In simulation, memistor-class three-terminal electrochemical weights can train small networks to the exact-update baseline,
provided the devices saturate, their reference is accurate (or a chopped algorithm removes the need for it) and a Tiki-Taka-family
algorithm moves the gradient. Two published specifications, Li-ECRAM and protonic, meet these conditions. Plain SGD does not
work on any of them. Under a fair per-arm tuning, the gap does not grow with width over the range we could test. Counting the
periphery, the energy advantage over a well-used H100 is about \REnRatioMin$\times$ today, rising past 100$\times$ only with read
periphery several times better than the best measured analog-AI chip. The remaining questions are experimental, and we have tried
to make them cheap to ask.
